\section{Vacancy arithmetic between the scales \texorpdfstring{\(3^6\)}{3 raised to 6} and \texorpdfstring{\(10^3\)}{10 raised to 3}}
\label{sec:vacancias-corriente-electron-rev11}
\HMTClaim{MD-R-VACANCY-POLARIZATION-CURRENT}
\index{vacancy!arithmetic calendar}
\index{polarization!discrete current}

Finite scales
\[
 729=3^6,
 \qquad
 1000=10^3
\]
determine an exact logarithmic mismatch. We define
\begin{equation}
 \rho_{\rm vac}=\log_{1000}729=\log_{10}9,
 \qquad
 \delta_{\rm vac}=1-\rho_{\rm vac}=\log_{10}(10/9).
 \label{eq:delta-vacancias-rev11}
\end{equation}
The irrationality of \(\delta_{\rm vac}\) follows directly from the prime valuations: an equality \(\delta_{\rm vac}=p/q\) would imply \((10/9)^q=10^p\), which is incompatible with the valuations at \(2,3,5\). The vacancy is therefore an aperiodic arithmetic rotation determined entirely by the relation between the two scales; its definition precedes the electronic law.

\subsection{Exact Laws of Hole and Return}

Let
\[
 p_n=\left\lfloor\frac{n}{\delta_{\rm vac}}\right\rfloor.
\]

\begin{theorem}[Gap law \texorpdfstring{\(21/22\)}{21/22}]
\label{thm:ley-huecos-vacancias-rev11}
The consecutive increments of the sequence \(p_n\) satisfy
\[
 p_{n+1}-p_n\in\{21,22\}.
\]
Both values appear in the word mechanical.
\end{theorem}

\begin{proof}
For every irrational \(x>1\), the difference
\(\lfloor(n+1)x\rfloor-\lfloor nx\rfloor\) belongs to
\(\{\lfloor x\rfloor,\lceil x\rceil\}\). In this case
\[
 \delta_{\rm vac}^{-1}
 =21.85434532678283256\ldots,
\]
which determines the two vacancies. Irrationality rules out stabilization
in either one alone.
\end{proof}

The frequency of the short hole is
\[
 \sigma_{\rm vac}=22-\delta_{\rm vac}^{-1}
 =0.1456546732171674374\ldots,
 \qquad
 \sigma_{\rm vac}^{-1}=6.865553832996660997\ldots.
\]

\begin{corollary}[Return law \texorpdfstring{\(6/7\)}{6/7}]
\label{cor:ley-retorno-vacancias-rev11}
Successive occurrences of the gap \(21\) are separated by six or
seven positions.
\end{corollary}

\begin{proof}
The indicator word for the short gap is the mechanical coding with
irrational slope \(\sigma_{\rm vac}\). Beatty's argument applied to
\(\sigma_{\rm vac}^{-1}\) produces exactly the returns \(6/7\)
\cite{Lothaire2002}.
\end{proof}

The two laws form the oriented framework
\begin{equation}
 V_{\rm vac}=
 \begin{pmatrix}21&22\\6&7\end{pmatrix},
 \qquad
 \det V_{\rm vac}=21\cdot7-22\cdot6=15.
 \label{eq:marco-vacancias-21-22-rev11}
\end{equation}
The nonzero determinant proves the independence of the gap and return
directions. The cardinal \(15\) thus reappears as an invariant of the
vacancy calendar. In the electron equation, the pair
\((-22,+15)\) comes from the central selector proved in
\cref{thm:coeficientes-electron-rev6}; the framework
\eqref{eq:marco-vacancias-21-22-rev11} provides an independent arithmetic
corroboration, with its own genealogy.

It is convenient to distinguish two words derived from the same mismatch. The
indicator
\[
 h_n=\mathbf 1_{\{p_{n+1}-p_n=21\}}
\]
has frequency \(\sigma_{\rm vac}\) and returns \(6/7\). The word
\[
 z_n=\lceil(n+1)\delta_{\rm vac}\rceil
      -\lceil n\delta_{\rm vac}\rceil
\]
has slope \(\delta_{\rm vac}\) and contains \(45\) or \(46\) units
in each block of length one thousand. The laws \(21/22\), \(6/7\), and \(45/46\)
are distinct readers of the same irrational rotation.

\subsection{Bounded polarization and discrete continuity equation}

A phase \(\theta\) defines
\[
 z_n(\theta)=\lceil(n+1)\delta_{\rm vac}+\theta\rceil
 -\lceil n\delta_{\rm vac}+\theta\rceil.
\]
The accumulated intrinsic polarization is
\begin{equation}
 \begin{aligned}
 P_N(\theta)
 &:=\sum_{n=0}^{N-1}\bigl(z_n(\theta)-\delta_{\rm vac}\bigr)\\
 &=\lceil N\delta_{\rm vac}+\theta\rceil-
   \lceil\theta\rceil-N\delta_{\rm vac},
 \qquad |P_N(\theta)|<1.
 \end{aligned}
 \label{eq:polarizacion-vacancias-rev11}
\end{equation}
The temporal difference
\begin{equation}
 J_N(\theta)=P_{N+1}(\theta)-P_N(\theta)
 =z_N(\theta)-\delta_{\rm vac}
 \label{eq:corriente-vacancias-rev11}
\end{equation}
is the associated arithmetic current. The
\cref{eq:polarizacion-vacancias-rev11,eq:corriente-vacancias-rev11}
constitute an exact continuity equation: the accumulated charge remains
uniformly bounded, and each pulse is its local variation. The electromagnetic
realization applies a unit of charge and a temporal scale to this pair;
the polarization and current structure is already fixed before that chart.

\subsection{Harmonic Closure of Conjugated Null Modes in the Electronic State}
\label{sec:nido-armonico-modos-nulos-rev11}
\HMTClaim{MD-R-ELECTRON-NULL-MODE-NEST-CANDIDATE}

The central electronic section admits a structural reading as the closure
of two conjugate null modes. For \(u,v>0\), the hyperbolic projectors
\(P_\pm\) satisfy
\[
 N_{\rm sp}(uP_++vP_-)=uv>0,
\]
although each isolated component has zero norm. In this construction, the resonant closure of conjugate electromagnetic null modes denotes the two-sheet transport ring whose holonomy closes the central material section; emission and absorption correspond to opening and closing a null orientation relative to that section.

\begin{statusbox}[title={Photonic realization of harmonic closure}]
The closure of two null modes and the positivity of their coupled norm are
algebraic identities. The electromagnetic realization is specified by
an intertwiner with the Maxwell modes, a closure Hamiltonian, and the
corresponding transition rates. The public expression
\emph{resonant closure of conjugate electromagnetic null modes}
names the structural object that receives that realization.
\end{statusbox}
